% ============================================================
% 纯答案册·课时10-2.4曲线与方程 —— 工具/答案抽册器.py 抽册生成件（勿手改；第三档＝纯答案单独成册）
% 源件（只读）：C:/提示词/工作区/M3-第2章量产0913/成卷/导学件/课时10-2.4曲线与方程/main.tex（md5 13cc4af64b704bad54e0ec41a76ca661）
%% 口径：题面不重印；逐题＝题号(印面号同正册)＋[答案]＋[详解]，值/详解照源件逐字
%       （钉值门硬断言：册值≡正册件内值≡值台账值tex，逐字全等）。册序＝正册装配序＝印面号 1..21 升序（同序同号）；键序断言基准＝题面库冻结 manifest canonical 键序。
%% 三档导出：整体 main-true／纯题 main-false（在产）｜本册＝纯答案册（本件）
%% 双档：ansbook-true.tex＝详解印本；ansbook-false.tex＝纯值速查本（详解不印）
% 编译：xelatex <壳> ×2，cwd＝本目录；sty＝qp-m3.sty 本地挂载（md5 7c3930362be8a0a2bdf21bbf8ac16573，锁校验）
% ============================================================
\documentclass[fontset=none]{ctexart}
\usepackage{qp-m3}
% —— 印面逐字符号路由补钉（承源件；− 走 TNR、▱ NSC 子块；禁改 sty 保 md5） ——
\xeCJKDeclareCharClass{Default}{"2212}%
\setCJKfamilyfont{nscblk}[Path=C:/提示词/工作区/字替对照-0909/variantF/fonts/,BoldFont=NSC-w600.ttf]{NSC-w500.ttf}%
\xeCJKDeclareSubCJKBlock{nscblk}{"25B1}%
\newcommand{\pxparallelogram}{{\CJKfamily{nscblk}▱}}%
% —— 断行微调（承母版实测档，三零门承重；\emergencystretch 不另设） ——
\xeCJKsetup{CJKglue={\hskip 0.10em plus 0.08em minus 0.05em}}%
% —— 纯答案册全线括线（长详解可跨栏断，承练习件全线括线口径；灰底盒不可跨栏断） ——
\ansblockgrayfalse
% —— 双档开关（false 档＝纯值速查：答案印、详解不印；件面开关，不动 sty 本体） ——
\newif\ifansbookdetail
\ifdefined\ansbookdetail
  \ifnum\ansbookdetail=0 \ansbookdetailfalse\else\ansbookdetailtrue\fi
\else
  \ansbookdetailtrue
\fi
% —— 页脚件名（承源件章名；件名＝<件型>答案册） ——
\renewcommand{\qpzhangming}{第二章\quad 平面解析几何}%
\renewcommand{\qpjianming}{导学件答案册}%

\begin{document}
\zhangtitle{第二章\quad 平面解析几何}
\jietitle{2.4 曲线与方程}
\par\glueguard{1}\addvspace{2pt}\noindent\biaoqian{【纯答案册】}{\fangsong 题面不重印\quad 印面号与正册同号同序}\par\addvspace{2pt}

\begin{multicols}{2}
\emergencystretch=1em

\begin{ansblock}[2章-练-课时10-简1]
% ans:2章-练-课时10-简1
\ansitem{1}{P在曲线上；Q不在曲线上}
\ifansbookdetail
\ansnote{详解}{曲线上的点的坐标必为方程的解，反之方程的解对应曲线上的点，故代入即判．\(P(\sqrt{3},0)\)：\(4\times(\sqrt{3})^{2}+3\times 0^{2}=12+0=12\)，是方程的解，\(P\)在曲线上．\(Q(-2,3)\)：\(4\times(-2)^{2}+3\times 3^{2}=16+27=43\neq 12\)，不是方程的解，\(Q\)不在曲线上．}
\fi
\end{ansblock}

\begin{ansblock}[2章-练-课时10-简2]
% ans:2章-练-课时10-简2
\ansitem{2}{r²=5}
\ifansbookdetail
\ansnote{详解}{曲线通过点\((-1,5)\)，即\((-1,5)\)的坐标是方程的解：\((-1+2)^{2}+(5-3)^{2}=r^{2}\)，得\(r^{2}=1+4=5\)．}
\fi
\end{ansblock}

\begin{ansblock}[2章-练-课时10-简3]
% ans:2章-练-课时10-简3
\ansitem{3}{不是。x−y=0 只表示其中一条直线，轨迹应为 x−y=0 或 x+y=0（即|x|=|y|）；以 x−y=0 为轨迹方程不完备（漏了直线 x+y=0 上的点）}
\ifansbookdetail
\ansnote{详解}{设动点\(M(x,y)\)．到\(x\)轴距离为\(|y|\)、到\(y\)轴距离为\(|x|\)，条件\(|x|=|y|\iff x-y=0\)或\(x+y=0\)，轨迹是两条互相垂直的直线．方程\(x-y=0\)的解对应的点都在轨迹上（纯粹性成立），但轨迹上的点（如\((1,-1)\)）不都是该方程的解（完备性不成立）——它只描述第一、三象限角平分线，漏了第二、四象限角平分线．故“轨迹的方程是\(x-y=0\)”不对．}
\fi
\end{ansblock}

\begin{ansblock}[2章-练-课时10-简7]
% ans:2章-练-课时10-简7
\ansitem{4}{A}
\ifansbookdetail
\ansnote{详解}{A：\(\sqrt{x^{2}}=|x|\)，两方程解集（图象）完全相同，表同一条曲线；B：\(y=x\)为直线，\(y=\sqrt{x^{2}}=|x|\)为折线（\(x<0\)时\(y=-x\)），异；C：\(|y|=x\iff x\geqslant 0\)且\(y=\pm x\)，为两条射线，而\(y=|x|\)还含\(x<0\)部分，异；D：\(x^{2}+y^{2}=0\iff\)点\((0,0)\)，\(y=0\)为整条\(x\)轴，异．故选：A．}
\fi
\end{ansblock}

\begin{ansblock}[2章-练-课时10-简8]
% ans:2章-练-课时10-简8
\ansitem{5}{C}
\ifansbookdetail
\ansnote{详解}{C 与原方程仅差移项，同解变形，解集恒等，表同一条曲线；A：\(\sqrt{4-x^{2}}\) 要求 \(y\geqslant 0\)，仅为上半圆，缺 \(y<0\) 部分不完备，排除；B：圆心 \((2,0)\)、半径 2 的圆，与圆心在原点的圆非同一条曲线，排除；D：\(|x|+|y|=2\) 是以 \((2,0)\)、\((0,2)\)、\((-2,0)\)、\((0,-2)\) 为顶点的正方形边界，点 \((1,1)\) 在其上却不在圆上，排除．故选：C．}
\fi
\end{ansblock}

\begin{ansblock}[2章-练-课时10-简10]
% ans:2章-练-课时10-简10
\ansitem{6}{(1) 假　(2) 真　(3) 真}
\ifansbookdetail
\ansnote{详解}{(1) \(x^{2}+y^{2}=0\) 当且仅当 \(x=y=0\)，表示原点一个点，非直线；(2) \(|x|=|y|\) 当且仅当 \(x^{2}=y^{2}\)（两侧均非负，平方同解），即 \(y^{2}=x^{2}\)，命题真；(3) 由曲线方程定义的完备性（以方程的解为坐标的点都在曲线上），命题真．}
\fi
\end{ansblock}

\begin{ansblock}[2章-练-课时10-简4]
% ans:2章-练-课时10-简4
\ansitem{7}{x+2y−7=0}
\ifansbookdetail
\ansnote{详解}{设\(P(x,y)\)为线段\(AB\)垂直平分线上任意一点，则\(|PA|=|PB|\)：\((x+1)^{2}+(y+1)^{2}=(x-3)^{2}+(y-7)^{2}\)，展开：\(x^{2}+2x+1+y^{2}+2y+1=x^{2}-6x+9+y^{2}-14y+49\)，化简得 \(8x+16y-56=0\)，即 \(x+2y-7=0\)．反向检验：满足 \(x+2y-7=0\) 的点均有\(|PA|=|PB|\)，在\(AB\)的垂直平分线上（两性完备）．验算：\(AB\)中点\((1,3)\)：\(1+6-7=0\)成立；\(k_{AB}=(7+1)/(3+1)=2\)，所求直线斜率\(-1/2\)，\(2\times(-1/2)=-1\)成立．}
\fi
\end{ansblock}

\begin{ansblock}[2章-练-课时10-简5]
% ans:2章-练-课时10-简5
\ansitem{8}{(x−7)²+(y−2)²=32}
\ifansbookdetail
\ansnote{详解}{设\(M(x,y)\)，\(|MA|/|MB|=\sqrt{2}\)，即\(|MA|^{2}=2|MB|^{2}\)：\((x+1)^{2}+(y-2)^{2}=2[(x-3)^{2}+(y-2)^{2}]\)．展开：\(x^{2}+2x+1+(y-2)^{2}=2x^{2}-12x+18+2(y-2)^{2}\)，移项整理：\(x^{2}-14x+(y-2)^{2}+17=0\)，配方得 \((x-7)^{2}+(y-2)^{2}=32\)．轨迹为以\((7,2)\)为圆心、\(4\sqrt{2}\)为半径的圆（阿波罗尼斯圆雏形，此名可不向学生出现）．验算：圆与直线\(AB\)（\(y=2\)）交于 \(M_1(7-4\sqrt{2},2)\)、\(M_2(7+4\sqrt{2},2)\)．对\(M_1\)：\(|M_1A|=|7-4\sqrt{2}+1|=8-4\sqrt{2}=4(2-\sqrt{2})\)，\(|M_1B|=|4\sqrt{2}-4|=4(\sqrt{2}-1)\)（取绝对值），比值\(=(2-\sqrt{2})/(\sqrt{2}-1)=\sqrt{2}(\sqrt{2}-1)/(\sqrt{2}-1)=\sqrt{2}\)成立；对\(M_2\)同法比值亦\(\sqrt{2}\)成立．}
\fi
\end{ansblock}

\begin{ansblock}[2章-练-课时10-简6]
% ans:2章-练-课时10-简6
\ansitem{9}{x²/81+y²/36=1（x≠0）}
\ifansbookdetail
\ansnote{详解}{设\(A(x,y)\)，\(x\neq 0\)（\(x=0\) 时 \(A\) 在 \(y\) 轴上与 \(B\)、\(C\) 共线，不能成三角形且斜率不存在）．\(k_{AB}\cdot k_{AC}=[(y-6)/x]\cdot[(y+6)/x]=(y^{2}-36)/x^{2}=-4/9\)，去分母：\(9(y^{2}-36)=-4x^{2}\)，整理得 \(4x^{2}+9y^{2}=324\)，即 \(x^{2}/81+y^{2}/36=1\)（\(x\neq 0\)）．}
\fi
\end{ansblock}

\begin{ansblock}[2章-练-课时10-简9]
% ans:2章-练-课时10-简9
\ansitem{10}{x²=8(y−2)}
\ifansbookdetail
\ansnote{详解}{设 \(M(x,y)\)．由题 \(|y|=\sqrt{x^{2}+(y-4)^{2}}\)．两边平方（两侧非负，等价）：\(y^{2}=x^{2}+y^{2}-8y+16\)，即 \(x^{2}=8(y-2)\)．检验：方程上任一点均满足定义（平方等价可逆），无杂点；\(A(0,4)\) 代入 \(16\neq 0\) 不在轨迹上，无需剔除．故轨迹方程为 \(x^{2}=8(y-2)\)．}
\fi
\end{ansblock}

\begin{ansblock}[2章-练-课时10-中2]
% ans:2章-练-课时10-中2
\ansitem{11}{x²/4+y²=1（x≠±2）}
\ifansbookdetail
\ansnote{详解}{设\(P(x,y)\)．斜率存在要求\(x\neq\pm 2\)；\(y^{2}/(x^{2}-4)=-1/4\Rightarrow 4y^{2}=4-x^{2}\)，即\(x^{2}/4+y^{2}=1\)（\(x\neq\pm 2\)）．纯粹性：轨迹上任一点的坐标推导步步可逆，都满足斜率方程（成立）；完备性：满足题设的点都满足\(x^{2}/4+y^{2}=1\)且\(x\neq\pm 2\)，都在轨迹上（成立）（\(y=0\)时斜率积为\(0\neq -1/4\)，端点自然排除）．故轨迹为椭圆\(x^{2}/4+y^{2}=1\)除去长轴两端点\((\pm 2,0)\)．}
\fi
\end{ansblock}

\begin{ansblock}[2章-练-课时10-中1]
% ans:2章-练-课时10-中1
\ansitem{12}{(1)两个外切于原点的单位圆（「8」字形）；(2)关于x轴、y轴、原点均对称；x∈[−2,2]，y∈[−1,1]；与y轴仅交于(0,0)，与x轴交于(±2,0)与(0,0)；S=2π}
\ifansbookdetail
\ansnote{详解}{(1)按\(|x|\)分段：\(x\geqslant 0\)时\((x-1)^{2}+y^{2}=1\)，\(x<0\)时\((x+1)^{2}+y^{2}=1\)，两圆圆心\((\pm 1,0)\)、半径1，外切于原点．(2)\(|x|\)使\(C\)关于\(y\)轴对称；两圆连心线在\(x\)轴上，故关于\(x\)轴对称；由两者得关于原点对称．范围：\(x\in[-2,2]\)，\(y\in[-1,1]\)．交点：\(x=0\Rightarrow y^{2}=0\)，仅\((0,0)\)；\(y=0\Rightarrow|x|=2\)或\(0\)，得\((\pm 2,0)\)与\((0,0)\)．两圆仅切于一点无重叠，\(S=2\times\pi\times 1^{2}=2\pi\)．}
\fi
\end{ansblock}

\begin{ansblock}[2章-练-课时10-中3]
% ans:2章-练-课时10-中3
\ansitem{13}{m∈[−6,−4]∪[0,2]}
\ifansbookdetail
\ansnote{详解}{圆心\(C_1(0,m)\)、\(C_2(0,-2)\)，半径\(r_1=3\)、\(r_2=1\)，圆心距\(|C_1C_2|=|m+2|\)．两圆有公共点当且仅当\(|r_1-r_2|\leqslant|C_1C_2|\leqslant r_1+r_2\)，即\(2\leqslant|m+2|\leqslant 4\)，解得\(2\leqslant m+2\leqslant 4\)或\(-4\leqslant m+2\leqslant -2\)，即\(m\in[-6,-4]\cup[0,2]\)．}
\fi
\end{ansblock}

\begin{ansblock}[2章-练-课时10-中4]
% ans:2章-练-课时10-中4
\ansitem{14}{k<9：椭圆；k=9：单点(1,2)；k>9：不表示任何曲线}
\ifansbookdetail
\ansnote{详解}{配方：\((x-1)^{2}+2(y-2)^{2}=9-k\)．\(k<9\)时右边为正，方程化为\((x-1)^{2}/(9-k)+(y-2)^{2}/\bigl((9-k)/2\bigr)=1\)，表示椭圆；\(k=9\)时方程为\((x-1)^{2}+2(y-2)^{2}=0\)，表示单点\((1,2)\)（退化椭圆）；\(k>9\)时左边非负而右边为负，无实数解，不表示任何曲线．}
\fi
\end{ansblock}

\begin{ansblock}[2章-练-课时10-难1]
% ans:2章-练-课时10-难1
\ansitem{15}{BCD}
\ifansbookdetail
\ansnote{详解}{A：令\(y=0\)得\(x^{4}=4x^{2}\)，\(x=0\)或\(\pm 2\)，整点\((2,0)\)、\((-2,0)\)、\((0,0)\)三个；若整点满足\(y\neq 0\)，则由B项结论\(x^{2}+y^{2}\leqslant 4\)只能\(y=\pm 1\)、\(x\in\{0,\pm 1,\pm 2\}\)：代入\((x^{2}+1)^{2}=4(x^{2}-1)\)，记\(X=x^{2}\)得\(X^{2}-2X+5=0\)，判别式\(4-20<0\)无解——\(y\neq 0\)无整点．故曲线仅过3个整点，“5个整点”错．B：由\((x^{2}+y^{2})^{2}=4(x^{2}-y^{2})\leqslant 4(x^{2}+y^{2})\)得\(x^{2}+y^{2}\leqslant 4\)，即\(|OP|\leqslant 2\)（成立）．C：方程中\(x\)、\(y\)互换得\((x^{2}+y^{2})^{2}=4(y^{2}-x^{2})\)，即关于\(y=x\)对称的曲线（成立）．D：\(y=kx\)代入：\(x^{4}(1+k^{2})^{2}=4x^{2}(1-k^{2})\)．\(x=0\)恒给交点\((0,0)\)；\(|x|>0\)的解须\(x^{2}=4(1-k^{2})/(1+k^{2})^{2}>0\)，当且仅当\(|k|<1\)．故\(|k|\geqslant 1\)时只有原点一个交点；\(|k|<1\)时另有\(\pm\)两交点（\((0,0)\)之外），不止一个（成立）（\(k=\pm 1\)时\(x^{4}\cdot 2=0\)仅原点，含入界点正确）．故选：BCD．}
\fi
\end{ansblock}

\begin{ansblock}[2章-练-课时10-难2]
% ans:2章-练-课时10-难2
\ansitem{16}{②④}
\ifansbookdetail
\ansnote{详解}{①\(m=1\)时\((x-1)^{2}+(y-2)^{2}=n^{2}/2\)，但\(n=0\)时退化为点\((1,2)\)，不（恒）表示圆——以偏概全，错误．②\(m=0\)、\(n=2\)：\(C\)：\(x^{2}+y^{2}=2\)．记\(D(3,3)\)，\(|DO|=3\sqrt{2}\)，切线长\(|DA|=|DB|=\sqrt{18-2}=4\)，切点\(A\)、\(B\)即\(C\)与圆\((x-3)^{2}+(y-3)^{2}=16\)（圆心\(D\)半径4）的公共点，公共弦方程由两圆方程相减：\(x^{2}+y^{2}-2-(x^{2}+y^{2}-6x-6y+2)=0\)，即\(3x+3y-2=0\)，正确．③\(m=1\)、\(n=\sqrt{2}\)：\(C\)：\((x-1)^{2}+(y-2)^{2}=1\)．过\((2,0)\)的切线：竖直直线\(x=2\)到圆心\((1,2)\)距离为\(1=\)半径，已是切线；另设\(y=k(x-2)\)：\(|k-2-2k|/\sqrt{k^{2}+1}=1\Rightarrow(k+2)^{2}=k^{2}+1\Rightarrow k=-3/4\)，即\(3x+4y-6=0\)也是切线．结论只给后者、漏\(x=2\)，“切线方程为”不完备，错误．④\(n=m\neq 0\)：\((x-m)^{2}+(y-2m)^{2}=m^{2}/2\)，圆心\((m,2m)\)恒在\(y=2x\)上，半径\((\sqrt{2}/2)|m|\)．圆心到\(y=x\)距离\(=|m-2m|/\sqrt{2}=(\sqrt{2}/2)|m|=\)半径；到\(y=7x\)距离\(=|7m-2m|/\sqrt{50}=(\sqrt{2}/2)|m|=\)半径——两直线都是圆系公切线．穷尽性补核（源未给）：设公切线\(y=kx+b\)对一切\(m\)成立，则\([m(k-2)+b]^{2}=\frac{1}{2}m^{2}(k^{2}+1)\)恒成立，比较系数：\(b=0\)且\((k-2)^{2}=(k^{2}+1)/2\Rightarrow k^{2}-8k+7=0\Rightarrow k=1\)或\(7\)；竖直线\(x=c\)代入\(|m-c|=(\sqrt{2}/2)|m|\)不能对一切\(m\)成立，无．故公切线恰为\(y=x\)、\(y=7x\)，正确．⑤\(n=4\)、\(m=0\)：\(C\)：\(x^{2}+y^{2}=8\)，直线\(kx-y+1-2k=0\)即\(k(x-2)-(y-1)=0\)恒过定点\((2,1)\)，而\(4+1=5<8\)，定点在圆内，直线恒与圆相交而非相离，错误．综上，正确的是②④．}
\fi
\end{ansblock}

\begin{ansblock}[2章-导-课时10-G1]
% ans:2章-导-课时10-G1
\ansitem{17}{D}
\ifansbookdetail
\ansnote{详解}{设\(M(x,y)\)，\(P(x_0,y_0)\)，则\(x_0^{2}+y_0^{2}=4\)．由\(PH\perp y\)轴知\(H(0,y_0)\)；\(M\)为\(PH\)中点，得\(x=x_0/2\)、\(y=y_0\)，即\(x_0=2x\)、\(y_0=y\)．代入圆方程：\(4x^{2}+y^{2}=4\)，整理得\(x^{2}+y^{2}/4=1\)．故选：D．}
\fi
\end{ansblock}

\begin{ansblock}[2章-导-课时10-G2]
% ans:2章-导-课时10-G2
\ansitem{18}{A}
\ifansbookdetail
\ansnote{详解}{\(F(-2,0)\)、\(C(2,0)\)，记\(F\)关于折痕\(l\)的对称点为\(A\)，\(l\)与\(AC\)交于点\(P\)，则\(A\)在圆周上，\(l\)为\(AF\)的垂直平分线，故\(|PA|=|PF|\)．于是\(|PF|+|PC|=|PA|+|PC|=|AC|=8>|FC|=4\)，点\(P\)的轨迹是以\(F\)、\(C\)为左、右焦点的椭圆，\(2a=8\)、\(2c=4\)，\(a=4\)、\(c=2\)、\(b=2\sqrt{3}\)，方程为\(x^{2}/16+y^{2}/12=1\)．故选：A．}
\fi
\end{ansblock}

\begin{ansblock}[2章-导-课时10-G3]
% ans:2章-导-课时10-G3
\ansitem{19}{D}
\ifansbookdetail
\ansnote{详解}{曲线定义域：\(x-1\geqslant 0\)且\(x^{2}+y^{2}-1>0\)．乘积为零当且仅当\(\sqrt{x-1}=0\)或\(\ln(x^{2}+y^{2}-1)=0\)．前者给\(x=1\)（此时\(x^{2}+y^{2}-1=y^{2}>0\)须\(y\neq 0\)），即直线\(x=1\)除去\((1,0)\)；后者给\(x^{2}+y^{2}=2\)且\(x\geqslant 1\)，即圆\(x^{2}+y^{2}=2\)的右半弧（含端点\((1,\pm 1)\)，该两点同时落在\(x=1\)上）．直线\(y=kx-1\)恒过定点\((0,-1)\)：与\(x=1\)（\(y\neq 0\)）交于\((1,k-1)\)，\(k\neq 1\)时是曲线上一交点；又\((0,-1)\)到原点距离\(1<\sqrt{2}\)，直线与整圆恒交于两点，其中落在右半弧上的：临界为直线过弧端点\((1,-1)\)（\(k=0\)）与\((1,1)\)（\(k=2\)），端点处弧上点与\(x=1\)上点重合、两交并一，故须\(0<k<2\)（严格）；再排除\(k=1\)（此时\((1,k-1)=(1,0)\)不在曲线上，曲线只剩弧上一个交点，\(k=1\)时直线与圆\(x\geqslant 1\)半弧恰一交点）．综上\(0<k<2\)且\(k\neq 1\)，故选：D．}
\fi
\end{ansblock}

\begin{ansblock}[2章-导-课时10-G4]
% ans:2章-导-课时10-G4
\ansitem{20}{x²+y²−8x−4y+18=0（除去点(3,1)与(5,3)）}
\ifansbookdetail
\ansnote{详解}{设\(C(x,y)\)．由\(|AC|=|AB|\)得\(\sqrt{(x-4)^{2}+(y-2)^{2}}=\sqrt{(4-5)^{2}+(2-3)^{2}}=\sqrt{2}\)，平方化简得\(x^{2}+y^{2}-8x-4y+18=0\)（圆心\((4,2)\)、半径\(\sqrt{2}\)的圆）．剔除：\(C\)与\(B\)重合，去\((5,3)\)；\(A\)、\(B\)、\(C\)三点共线时不构成三角形——\(k_{AB}=(3-2)/(5-4)=1\)，直线\(AB\)：\(y-2=1\cdot(x-4)\)即\(x-y-2=0\)，与圆联立解得交点\((3,1)\)与\((5,3)\)，故须除去\((3,1)\)、\((5,3)\)．轨迹方程为\(x^{2}+y^{2}-8x-4y+18=0\)（\(x-y-2\neq 0\)，即除去\((3,1)\)、\((5,3)\)）．}
\fi
\end{ansblock}

\begin{ansblock}[2章-导-课时10-G5]
% ans:2章-导-课时10-G5
\ansitem{21}{x=4 或 x=−4（即 x=±4，两条平行于y轴的直线）}
\ifansbookdetail
\ansnote{详解}{设动点\(M(x,y)\)．\(M\)到\(y\)轴距离为\(|x|\)，条件\(|x|=4\)即\(x=4\)或\(x=-4\)．检验两性：直线\(x=\pm 4\)上任意一点到\(y\)轴距离均为4（纯粹性成立）；到\(y\)轴距离为4的点横坐标为\(\pm 4\)，都在两直线上（完备性成立）．故轨迹方程为\(x=\pm 4\)．}
\fi
\end{ansblock}

% ---- 尾块（\tailfill[30mm] 定高参——钉档 §三.3：无参在平衡栏呈「内容尾随框」，贴栏底须定高参；实测无参致末页平衡 Overfull\vbox 12.99pt，定高 30mm 三零） ----
\tailfill[30mm]

\end{multicols}
\end{document}
